\section{Discussion}
\label{sec:discussion}

\para{What do the results show?} The main finding is not that LLMs can fingerprint arbitrary generators. The more defensible claim is narrower: current LLMs can read broad production modes from text alone, and that signal is strong enough to appear both in direct judgments and in embedding-space analyses. The controlled benchmark, grouped probe, and delta-consistency tests all support that view. The external \semeval result then provides the necessary boundary condition by showing that zero-shot generator identity is much less stable.

\para{Why does the split matter?} If a model can separate dictated text, keyboard-noisy text, and LLM-post-edited prose, then mode already affects systems that use LLMs as judges or filters. Provenance screening, moderation triage, and educational-integrity workflows may therefore be reacting to more than semantic content. At the same time, our external result argues against overinterpreting these abilities as reliable model fingerprinting. In practice, a detector that is good at reading process state may still be poor at naming the exact generator family.

\para{The latent-variable evidence is promising but partial.} The grouped probe and the delta-similarity tests support the idea that mode behaves like a reusable direction in representation space. That matters because it suggests mode is not merely a prompt artifact attached to one benchmark. Still, the evidence is not yet definitive. The perfect custom classification score indicates that the benchmark is relatively easy, and same-family generation and classification may amplify the signal.

\para{Limitations.} Several threats to validity remain. First, \gptfourone generated several transformed modes and also served as the judge, which raises a self-recognition risk. Second, the \texttt{keyboard\_noisy} mode is synthetic and deterministic rather than collected from real typing data. Third, content is only approximately matched across variants, so some semantic drift may remain. Fourth, the permutation test is exploratory because it uses only 20 corrected runs. Fifth, the \semeval experiment is zero-shot only; a tuned attribution model would almost certainly perform much better. Finally, our analysis is document-level and does not test span-level mode transitions such as those in \semeval Subtask C.

\para{Broader implications.} These results support a more cautious vocabulary for provenance claims. Broad production mode is likely a real and useful construct for ambiguity-aware evaluation and screening, but exact source attribution should be treated as a separate and harder problem. We expect the strongest near-term applications to involve abstaining or flagging mode uncertainty rather than making high-confidence claims about exact model identity.
